\documentclass[confidential]{nervtech-engdoc}

\nervsetup{
  company={NervTech},
  project={Integrated Robotic Actuator},
  document={Integrated Robotic Actuator Technical Design Document},
  subtitle={Electromechanical architecture, control design, verification, and integration plan},
  docid={NT-IRA-TDD-001},
  revision={A},
  status={Draft for Design Review},
  author={NervTech Engineering},
  owner={Robotics Systems Engineering},
  approver={Project Lead},
  date={2026-06-25},
  keywords={robotic actuator, BLDC, embedded control, torque sensing, mechatronics},
  abstract={This document captures the architecture, requirements, safety constraints, implementation notes, and verification plan for the NervTech Integrated Robotic Actuator project.}
}

\begin{document}
\maketitle
\tableofcontents
\newpage

\begin{revisionhistory}
\revisionrow{A}{2026-06-25}{NervTech Engineering}{Initial technical document class demonstration.}
\end{revisionhistory}

\section{System Overview}
\systemsummary
  {Compact BLDC actuator with gearbox, encoder feedback, and embedded motor controller.}
  {Cascaded current, velocity, and position control with torque-limit supervision.}
  {Motor phase current, rotor position, output position, temperature, and bus voltage.}
  {Fused DC input, firmware watchdog, current limit, thermal derating, and emergency stop path.}

\begin{designnote}[Architecture Intent]
The integrated actuator should be documented as a full mechatronic system: motor, transmission, bearings, enclosure, PCB, firmware, control loops, and validation fixtures.
\end{designnote}

\section{Requirements}
\begin{requirementstable}
\reqrow{IRA-REQ-001}{Functional}{The actuator shall support closed-loop position control using encoder feedback.}{High}{Bench test}
\reqrow{IRA-REQ-002}{Safety}{The actuator shall enter a disabled state if bus voltage, phase current, or temperature exceeds configured limits.}{Critical}{Fault injection}
\reqrow{IRA-REQ-003}{Interface}{The actuator shall expose telemetry over a deterministic digital interface.}{Medium}{Integration test}
\end{requirementstable}

\section{Electrical Architecture}
\begin{safetynote}[Power Bring-up]
Start with a current-limited bench supply, no load, conservative firmware limits, and an accessible emergency disconnect before energising the motor stage.
\end{safetynote}

\begin{interfacetable}
\interfacerow{DC Bus}{Power}{24--48 V input, fused}{Ripple and inrush must remain within controller limits}
\interfacerow{Encoder}{Signal}{Quadrature or SPI absolute position}{Shielding and grounding required near motor phases}
\interfacerow{Host Link}{Data}{CAN, RS-485, or UART telemetry}{Deterministic update period for control logging}
\end{interfacetable}

\section{Firmware Example}
\begin{lstlisting}[language=C++,caption={Torque-limit supervision pseudocode}]
if (phase_current > current_limit || motor_temp > temp_limit) {
    controller.disable_pwm();
    fault_latch = true;
    publish_fault_status();
}
\end{lstlisting}

\section{Verification Plan}
\begin{testcasetable}
\testrow{IRA-TST-001}{Command a 10 degree step response at no load.}{Settling time and overshoot are within design target.}{TBD}
\testrow{IRA-TST-002}{Inject over-current threshold in firmware test mode.}{PWM is disabled and fault is latched.}{TBD}
\end{testcasetable}

\end{document}
